% Options for packages loaded elsewhere
% Options for packages loaded elsewhere
\PassOptionsToPackage{unicode}{hyperref}
\PassOptionsToPackage{hyphens}{url}
\PassOptionsToPackage{dvipsnames,svgnames,x11names}{xcolor}
%
\documentclass[
  letterpaper,
  DIV=11,
  numbers=noendperiod]{scrartcl}
\usepackage{xcolor}
\usepackage{amsmath,amssymb}
\setcounter{secnumdepth}{5}
\usepackage{iftex}
\ifPDFTeX
  \usepackage[T1]{fontenc}
  \usepackage[utf8]{inputenc}
  \usepackage{textcomp} % provide euro and other symbols
\else % if luatex or xetex
  \usepackage{unicode-math} % this also loads fontspec
  \defaultfontfeatures{Scale=MatchLowercase}
  \defaultfontfeatures[\rmfamily]{Ligatures=TeX,Scale=1}
\fi
\usepackage{lmodern}
\ifPDFTeX\else
  % xetex/luatex font selection
\fi
% Use upquote if available, for straight quotes in verbatim environments
\IfFileExists{upquote.sty}{\usepackage{upquote}}{}
\IfFileExists{microtype.sty}{% use microtype if available
  \usepackage[]{microtype}
  \UseMicrotypeSet[protrusion]{basicmath} % disable protrusion for tt fonts
}{}
% Make \paragraph and \subparagraph free-standing
\makeatletter
\ifx\paragraph\undefined\else
  \let\oldparagraph\paragraph
  \renewcommand{\paragraph}{
    \@ifstar
      \xxxParagraphStar
      \xxxParagraphNoStar
  }
  \newcommand{\xxxParagraphStar}[1]{\oldparagraph*{#1}\mbox{}}
  \newcommand{\xxxParagraphNoStar}[1]{\oldparagraph{#1}\mbox{}}
\fi
\ifx\subparagraph\undefined\else
  \let\oldsubparagraph\subparagraph
  \renewcommand{\subparagraph}{
    \@ifstar
      \xxxSubParagraphStar
      \xxxSubParagraphNoStar
  }
  \newcommand{\xxxSubParagraphStar}[1]{\oldsubparagraph*{#1}\mbox{}}
  \newcommand{\xxxSubParagraphNoStar}[1]{\oldsubparagraph{#1}\mbox{}}
\fi
\makeatother

\usepackage{color}
\usepackage{fancyvrb}
\newcommand{\VerbBar}{|}
\newcommand{\VERB}{\Verb[commandchars=\\\{\}]}
\DefineVerbatimEnvironment{Highlighting}{Verbatim}{commandchars=\\\{\}}
% Add ',fontsize=\small' for more characters per line
\usepackage{framed}
\definecolor{shadecolor}{RGB}{241,243,245}
\newenvironment{Shaded}{\begin{snugshade}}{\end{snugshade}}
\newcommand{\AlertTok}[1]{\textcolor[rgb]{0.68,0.00,0.00}{#1}}
\newcommand{\AnnotationTok}[1]{\textcolor[rgb]{0.37,0.37,0.37}{#1}}
\newcommand{\AttributeTok}[1]{\textcolor[rgb]{0.40,0.45,0.13}{#1}}
\newcommand{\BaseNTok}[1]{\textcolor[rgb]{0.68,0.00,0.00}{#1}}
\newcommand{\BuiltInTok}[1]{\textcolor[rgb]{0.00,0.23,0.31}{#1}}
\newcommand{\CharTok}[1]{\textcolor[rgb]{0.13,0.47,0.30}{#1}}
\newcommand{\CommentTok}[1]{\textcolor[rgb]{0.37,0.37,0.37}{#1}}
\newcommand{\CommentVarTok}[1]{\textcolor[rgb]{0.37,0.37,0.37}{\textit{#1}}}
\newcommand{\ConstantTok}[1]{\textcolor[rgb]{0.56,0.35,0.01}{#1}}
\newcommand{\ControlFlowTok}[1]{\textcolor[rgb]{0.00,0.23,0.31}{\textbf{#1}}}
\newcommand{\DataTypeTok}[1]{\textcolor[rgb]{0.68,0.00,0.00}{#1}}
\newcommand{\DecValTok}[1]{\textcolor[rgb]{0.68,0.00,0.00}{#1}}
\newcommand{\DocumentationTok}[1]{\textcolor[rgb]{0.37,0.37,0.37}{\textit{#1}}}
\newcommand{\ErrorTok}[1]{\textcolor[rgb]{0.68,0.00,0.00}{#1}}
\newcommand{\ExtensionTok}[1]{\textcolor[rgb]{0.00,0.23,0.31}{#1}}
\newcommand{\FloatTok}[1]{\textcolor[rgb]{0.68,0.00,0.00}{#1}}
\newcommand{\FunctionTok}[1]{\textcolor[rgb]{0.28,0.35,0.67}{#1}}
\newcommand{\ImportTok}[1]{\textcolor[rgb]{0.00,0.46,0.62}{#1}}
\newcommand{\InformationTok}[1]{\textcolor[rgb]{0.37,0.37,0.37}{#1}}
\newcommand{\KeywordTok}[1]{\textcolor[rgb]{0.00,0.23,0.31}{\textbf{#1}}}
\newcommand{\NormalTok}[1]{\textcolor[rgb]{0.00,0.23,0.31}{#1}}
\newcommand{\OperatorTok}[1]{\textcolor[rgb]{0.37,0.37,0.37}{#1}}
\newcommand{\OtherTok}[1]{\textcolor[rgb]{0.00,0.23,0.31}{#1}}
\newcommand{\PreprocessorTok}[1]{\textcolor[rgb]{0.68,0.00,0.00}{#1}}
\newcommand{\RegionMarkerTok}[1]{\textcolor[rgb]{0.00,0.23,0.31}{#1}}
\newcommand{\SpecialCharTok}[1]{\textcolor[rgb]{0.37,0.37,0.37}{#1}}
\newcommand{\SpecialStringTok}[1]{\textcolor[rgb]{0.13,0.47,0.30}{#1}}
\newcommand{\StringTok}[1]{\textcolor[rgb]{0.13,0.47,0.30}{#1}}
\newcommand{\VariableTok}[1]{\textcolor[rgb]{0.07,0.07,0.07}{#1}}
\newcommand{\VerbatimStringTok}[1]{\textcolor[rgb]{0.13,0.47,0.30}{#1}}
\newcommand{\WarningTok}[1]{\textcolor[rgb]{0.37,0.37,0.37}{\textit{#1}}}

\usepackage{longtable,booktabs,array}
\usepackage{calc} % for calculating minipage widths
% Correct order of tables after \paragraph or \subparagraph
\usepackage{etoolbox}
\makeatletter
\patchcmd\longtable{\par}{\if@noskipsec\mbox{}\fi\par}{}{}
\makeatother
% Allow footnotes in longtable head/foot
\IfFileExists{footnotehyper.sty}{\usepackage{footnotehyper}}{\usepackage{footnote}}
\makesavenoteenv{longtable}
\usepackage{graphicx}
\makeatletter
\newsavebox\pandoc@box
\newcommand*\pandocbounded[1]{% scales image to fit in text height/width
  \sbox\pandoc@box{#1}%
  \Gscale@div\@tempa{\textheight}{\dimexpr\ht\pandoc@box+\dp\pandoc@box\relax}%
  \Gscale@div\@tempb{\linewidth}{\wd\pandoc@box}%
  \ifdim\@tempb\p@<\@tempa\p@\let\@tempa\@tempb\fi% select the smaller of both
  \ifdim\@tempa\p@<\p@\scalebox{\@tempa}{\usebox\pandoc@box}%
  \else\usebox{\pandoc@box}%
  \fi%
}
% Set default figure placement to htbp
\def\fps@figure{htbp}
\makeatother





\setlength{\emergencystretch}{3em} % prevent overfull lines

\providecommand{\tightlist}{%
  \setlength{\itemsep}{0pt}\setlength{\parskip}{0pt}}



 


\usepackage{booktabs}
\usepackage{caption}
\usepackage{longtable}
\usepackage{colortbl}
\usepackage{array}
\usepackage{anyfontsize}
\usepackage{multirow}
\usepackage{xcolor, soul}
\definecolor{Light}{gray}{.90}
\sethlcolor{Light}
\let\OldTexttt\texttt
\renewcommand{\texttt}[1]{\OldTexttt{\hl{#1}}}
\usepackage{caption}
\captionsetup{font=small, labelfont=bf}
\setlength{\parindent}{0pt}
\setlength{\parskip}{1.25ex}
\KOMAoption{captions}{tableheading}
\makeatletter
\@ifpackageloaded{caption}{}{\usepackage{caption}}
\AtBeginDocument{%
\ifdefined\contentsname
  \renewcommand*\contentsname{Table of contents}
\else
  \newcommand\contentsname{Table of contents}
\fi
\ifdefined\listfigurename
  \renewcommand*\listfigurename{List of Figures}
\else
  \newcommand\listfigurename{List of Figures}
\fi
\ifdefined\listtablename
  \renewcommand*\listtablename{List of Tables}
\else
  \newcommand\listtablename{List of Tables}
\fi
\ifdefined\figurename
  \renewcommand*\figurename{Figure}
\else
  \newcommand\figurename{Figure}
\fi
\ifdefined\tablename
  \renewcommand*\tablename{Table}
\else
  \newcommand\tablename{Table}
\fi
}
\@ifpackageloaded{float}{}{\usepackage{float}}
\floatstyle{ruled}
\@ifundefined{c@chapter}{\newfloat{codelisting}{h}{lop}}{\newfloat{codelisting}{h}{lop}[chapter]}
\floatname{codelisting}{Listing}
\newcommand*\listoflistings{\listof{codelisting}{List of Listings}}
\makeatother
\makeatletter
\makeatother
\makeatletter
\@ifpackageloaded{caption}{}{\usepackage{caption}}
\@ifpackageloaded{subcaption}{}{\usepackage{subcaption}}
\makeatother
\usepackage{bookmark}
\IfFileExists{xurl.sty}{\usepackage{xurl}}{} % add URL line breaks if available
\urlstyle{same}
\hypersetup{
  pdftitle={Loading stimuli labels and MUA data from MATLAB files},
  pdfauthor={Al Xin},
  colorlinks=true,
  linkcolor={blue},
  filecolor={Maroon},
  citecolor={Blue},
  urlcolor={Blue},
  pdfcreator={LaTeX via pandoc}}


\title{Loading stimuli labels and MUA data from MATLAB files}
\author{Al Xin}
\date{May 21, 2026}
\begin{document}
\maketitle


\section{Stimulus labels}\label{stimulus-labels}

The table of stimuli information contains the family ID, coherence,
sample ID, and block placement. Each unique combination of family ID,
coherence, and sample ID corresponds to a specific image, and each image
is repeated four times in the session. Blank images are arbitrarily
labeled with family ID \texttt{0} and sample ID \texttt{id\_blank} from
1 to 100.

Family ID, coherence, and sample ID columns are identical across all
files and only need to be saved once. However, there is a column marking
the position of the stimulus within the block, which is unique for each
session file. I cannot find the order of the blocks in the published
data of the blocks. This setup does not generate unique rows because the
repetition number of each unique image (defined by the combination of
coherence, family, and sample ID) is not recorded in any column and
block placement is arbitrary and may be repeated. I assume the
repetition number was implicitly captured by the row number, but I am
not confident this is the scheme Lee et al.~intended.

\subsection{Column details}\label{column-details}

\begin{itemize}
\tightlist
\item
  \texttt{coherence}: \texttt{0} when blank, \texttt{1} when noisy,
  \texttt{2} when naturalistic.
\item
  \texttt{family}: \texttt{0} when blank, two or three digit integer
  code otherwise.
\item
  \texttt{id}: \texttt{0} when blank, integer from \texttt{1} to
  \texttt{15} otherwise.
\item
  \texttt{id\_blank}: integer from \texttt{1} to \texttt{100} if blank,
  \texttt{0} otherwise.
\item
  \texttt{repetition}: column constructed from counting the repetition
  of \texttt{coherence}, \texttt{family}, and \texttt{id} or
  \texttt{id\_blank} combo.
\end{itemize}

\begin{Shaded}
\begin{Highlighting}[]
\CommentTok{\# get arbitrary identical stimulus labels}
\NormalTok{stims\_ref }\OtherTok{\textless{}{-}} \FunctionTok{list.files}\NormalTok{(t2\_dir) }\SpecialCharTok{\%\textgreater{}\%}
\NormalTok{  .[}\DecValTok{1}\NormalTok{] }\SpecialCharTok{\%\textgreater{}\%} \CommentTok{\# arbitrary index}
  \FunctionTok{paste0}\NormalTok{(t2\_dir, .) }\SpecialCharTok{\%\textgreater{}\%}
\NormalTok{  R.matlab}\SpecialCharTok{::}\FunctionTok{readMat}\NormalTok{() }\SpecialCharTok{\%\textgreater{}\%}
\NormalTok{  .}\SpecialCharTok{$}\NormalTok{save.label }\SpecialCharTok{\%\textgreater{}\%}
  \FunctionTok{data.frame}\NormalTok{()}
\FunctionTok{colnames}\NormalTok{(stims\_ref) }\OtherTok{\textless{}{-}} \FunctionTok{c}\NormalTok{(}
  \StringTok{"coherence"}\NormalTok{, }\StringTok{"family"}\NormalTok{, }\StringTok{"id"}\NormalTok{, }\StringTok{"coherence\_t3"}\NormalTok{, }\StringTok{"x5"}\NormalTok{, }\StringTok{"id\_blank"}\NormalTok{)}
\NormalTok{stimuli }\OtherTok{\textless{}{-}} \FunctionTok{select}\NormalTok{(stims\_ref, }\SpecialCharTok{{-}}\NormalTok{coherence\_t3, }\SpecialCharTok{{-}}\NormalTok{x5) }\SpecialCharTok{\%\textgreater{}\%}
  \FunctionTok{group\_by}\NormalTok{(coherence, family, id, id\_blank) }\SpecialCharTok{\%\textgreater{}\%}
  \FunctionTok{mutate}\NormalTok{(}\AttributeTok{repetition =} \FunctionTok{row\_number}\NormalTok{()) }\SpecialCharTok{\%\textgreater{}\%}
  \FunctionTok{ungroup}\NormalTok{()}
\end{Highlighting}
\end{Shaded}

\section{Binned multiunit activity}\label{binned-multiunit-activity}

Reading from MATLAB tends to be time-consuming, so it is easier to work
with the binned multiunit activity (MUA) pre-extracted and saved in R
data files with the quick serialization package \texttt{qs2}. The
locations of \texttt{NA} values were removed. I manually recoded the
subject IDs to be consistent with the monkey IDs in the table and
arbitrarily labeled the arrays \texttt{a} or \texttt{b} with the special
cases of \texttt{a} being shared for the V1/V2 border array on M1 and
\texttt{c} used for the second implantation of M6
(Table~\ref{tbl-subjects}). This separated the 17 original session files
into 27 subject files.

Furthermore, though the published MATLAB files contain the subject IDs,
there does not appear to be a mapping of the binned data to the subject.
Again, I assume that the mapping is implicitly captured by the ordering
and split the binned data into ordered chunks of 96 sites.

\begin{Shaded}
\begin{Highlighting}[]
\FunctionTok{data.frame}\NormalTok{(}\AttributeTok{ada\_id =}\NormalTok{ id\_recoding) }\SpecialCharTok{\%\textgreater{}\%}
\NormalTok{  tibble}\SpecialCharTok{::}\FunctionTok{rownames\_to\_column}\NormalTok{(}\StringTok{"lee\_id"}\NormalTok{) }\SpecialCharTok{\%\textgreater{}\%}
  \FunctionTok{relocate}\NormalTok{(lee\_id, }\AttributeTok{.before =}\NormalTok{ ada\_id) }\SpecialCharTok{\%\textgreater{}\%}
  \FunctionTok{gt}\NormalTok{() }\SpecialCharTok{\%\textgreater{}\%}
  \FunctionTok{tab\_options}\NormalTok{(}\AttributeTok{table.font.size =} \FunctionTok{px}\NormalTok{(}\DecValTok{12}\NormalTok{)) }\SpecialCharTok{\%\textgreater{}\%} 
  \FunctionTok{cols\_label}\NormalTok{(}
    \AttributeTok{lee\_id =} \StringTok{"Lee et al. IDs"}\NormalTok{, }
    \AttributeTok{ada\_id =} \StringTok{"IDs in ADA"}\NormalTok{)}
\end{Highlighting}
\end{Shaded}

\begin{table}

\caption{\label{tbl-subjects}Original MATLAB subject IDs from Lee et
al.~(left) and new IDs for this project (right).}

\centering{

\fontsize{9.0pt}{11.0pt}\selectfont
\begin{tabular*}{\linewidth}{@{\extracolsep{\fill}}ll}
\toprule
Lee et al. IDs & IDs in ADA \\ 
\midrule\addlinespace[2.5pt]
KY\_n1 & M1a \\ 
KY\_v1 & M1a \\ 
KY\_n2 & M1b \\ 
LI\_n1 & M2a \\ 
LI\_n2 & M2b \\ 
BB\_n1 & M3a \\ 
BB\_n2 & M3b \\ 
HT\_V4 & M6a \\ 
HT\_IT & M6b \\ 
HT\_n2 & M6c \\ 
HU\_n2 & M7a \\ 
XT\_n2 & M8a \\ 
\bottomrule
\end{tabular*}

}

\end{table}%

\begin{Shaded}
\begin{Highlighting}[]
\NormalTok{n }\OtherTok{\textless{}{-}} \DecValTok{96} \CommentTok{\# number of channels}
\NormalTok{files }\OtherTok{\textless{}{-}} \FunctionTok{list.files}\NormalTok{(t2\_dir)}

\ControlFlowTok{for}\NormalTok{ (f }\ControlFlowTok{in}\NormalTok{ files) \{}
\NormalTok{  f\_details }\OtherTok{\textless{}{-}} \FunctionTok{file\_info}\NormalTok{(f)}
\NormalTok{  dat }\OtherTok{\textless{}{-}} \FunctionTok{paste0}\NormalTok{(t2\_dir, }\StringTok{"/"}\NormalTok{, f) }\SpecialCharTok{\%\textgreater{}\%}\NormalTok{ R.matlab}\SpecialCharTok{::}\FunctionTok{readMat}\NormalTok{()}
\NormalTok{  subj }\OtherTok{\textless{}{-}} \FunctionTok{as.character}\NormalTok{(dat}\SpecialCharTok{$}\NormalTok{incl.subj)}
\NormalTok{  bins }\OtherTok{\textless{}{-}}\NormalTok{ dat}\SpecialCharTok{$}\NormalTok{bins}
  
  \CommentTok{\# start messaging block}
  \FunctionTok{cli\_text}\NormalTok{(}\FunctionTok{paste0}\NormalTok{(}
    \StringTok{"Session \{f\_details$region\} at \{f\_details$age\} weeks"}\NormalTok{, }\StringTok{" "}\NormalTok{, }
    \StringTok{"(\{f\_details$n\_subjects\} subjects)"}\NormalTok{))}
\NormalTok{  ulid }\OtherTok{\textless{}{-}} \FunctionTok{cli\_ul}\NormalTok{()}
  
  \CommentTok{\# walk through subjects in file}
  \ControlFlowTok{for}\NormalTok{ (i }\ControlFlowTok{in} \DecValTok{1}\SpecialCharTok{:}\FunctionTok{length}\NormalTok{(subj)) \{}
    \CommentTok{\# get the more informative subject code}
\NormalTok{    id\_code }\OtherTok{\textless{}{-}}\NormalTok{ id\_recoding[subj[i]]}
    \CommentTok{\# incorporate into file name (qs2 extension)}
\NormalTok{    f\_name }\OtherTok{\textless{}{-}} \FunctionTok{paste0}\NormalTok{(}
\NormalTok{      f\_details}\SpecialCharTok{$}\NormalTok{region, }\StringTok{"\_"}\NormalTok{, }
\NormalTok{      f\_details}\SpecialCharTok{$}\NormalTok{age, }\StringTok{"\_"}\NormalTok{, }
\NormalTok{      id\_code, }\StringTok{".qs2"}\NormalTok{)}
    \FunctionTok{cli\_li}\NormalTok{(}\StringTok{"\{i\}/\{length(subj)\}: saving \{subj[i]\} to \{f\_name\}..."}\NormalTok{)}
    \CommentTok{\# message(i, "/", length(subj),  ": ", subj[i], " to ", f\_name, "...")}
    
    \CommentTok{\# manually subset the MUA data by number of channels}
\NormalTok{    idx\_last }\OtherTok{\textless{}{-}}\NormalTok{ i }\SpecialCharTok{*}\NormalTok{ n}
\NormalTok{    idx\_first }\OtherTok{\textless{}{-}}\NormalTok{ idx\_last }\SpecialCharTok{{-}}\NormalTok{ n }\SpecialCharTok{+} \DecValTok{1}
\NormalTok{    bins\_ }\OtherTok{\textless{}{-}}\NormalTok{ bins[, , idx\_first}\SpecialCharTok{:}\NormalTok{idx\_last]}
    
    \CommentTok{\# detect missing channels in 3rd dimension}
\NormalTok{    missing\_sites }\OtherTok{\textless{}{-}} \FunctionTok{is.na}\NormalTok{(bins\_[}\DecValTok{1}\NormalTok{, }\DecValTok{1}\NormalTok{, ])}
\NormalTok{    bins\_ }\OtherTok{\textless{}{-}}\NormalTok{ bins\_[, , }\SpecialCharTok{!}\NormalTok{missing\_sites]}
    \ControlFlowTok{if}\NormalTok{ (}\FunctionTok{sum}\NormalTok{(missing\_sites) }\SpecialCharTok{\textgreater{}} \DecValTok{0}\NormalTok{)}
      \CommentTok{\# message("NOTE: ", sum(missing\_sites), " missing channels.")}
      \FunctionTok{cli\_alert\_info}\NormalTok{(}\StringTok{"Found \{sum(missing\_sites)\} missing sites."}\NormalTok{)}
    
    \FunctionTok{qs\_save}\NormalTok{(bins\_, }\FunctionTok{paste0}\NormalTok{(mua\_dir, f\_name))}
\NormalTok{  \}}
  \CommentTok{\# close messaging block}
  \FunctionTok{cli\_end}\NormalTok{(ulid)}
  \FunctionTok{cli\_rule}\NormalTok{()}
\NormalTok{\}}
\end{Highlighting}
\end{Shaded}





\end{document}
